\documentclass[main]{subfiles}
\usepackage[sols]{workbook}
\numbering{ws2-8}

\begin{document}

\ifSubfilesClassLoaded{%
  \chaptermark{\chaptertwo}
}{}

\section{The Asymptotic Comparison Test} \label{ws2-8}
\begin{framed}
\centerline{\textbf{Lesson Goals}}

You should be able to:
\begin{itemize}
\item Determine if the terms from two series are asymptotic to each other using limits
\item Use growth rates to asymptotically simplify the terms of a series, and decide when an expression cannot be simplified further
\item Apply the Asymptotic Comparison Test to determine if a given series converges
\item Apply a combination of comparison strategies to determine if a given series converges

\end{itemize}
\end{framed} 


\begin{enumerate}
    \item 
    \note{I use this to introduce the ``problem'' --- here the series
    ``behaves like'' $\sum \frac{3}{2^k}$, but direct comparison doesn't
    work.  We'll define $\sim$ and state the Asymptotic Comparison Test
    and then use that to do this problem.  (When I do this in class, I
    will use the limit definition of $\sim$ to explain why $2^k - 100 \sim
    2^k$; the solutions use a different method.)}

    \begin{problem}[1.5in]
    Does the series $\displaystyle \sum_{k=1}^\infty \frac{3}{2^k - 100}$
    converge or diverge?
    \end{problem}

    \begin{solution}
    Since $2^k \gg 100$, $2^k - 100 \sim 2^k$, so $\displaystyle
    \frac{3}{2^k - 100} \sim \frac{3}{2^k}$.  Since $\displaystyle
    \frac{3}{2^k} > 0$, we can apply the Asymptotic Comparison Test to
    conclude that $\displaystyle \sum \frac{3}{2^k - 100}$ behaves the
    same way as $\displaystyle \sum \frac{3}{2^k}$.  The series
    $\displaystyle \sum \frac{3}{2^k}$ is geometric with common ratio
    $\displaystyle \frac{1}{2}$, so it converges.  Therefore,
    $\displaystyle \sum_{k=1}^\infty \frac{3}{2^k - 100}$ converges as
    well.
    \end{solution}
  \end{enumerate}
  
  
  \probonly{\vfill}
  

  \begin{framed}
  \textbf{Definition.}  Suppose $a_k$ and $b_k$ represent positive terms (eventually). We say that $b_k$ \uline{is asymptotically larger than} $a_k$
  as $k \to \infty$ if $\boxed{\displaystyle \lim_{k \to \infty}
  \frac{a_k}{b_k} = 0}$ or, equivalently, $\boxed{\displaystyle \lim_{k \to
  \infty} \frac{b_k}{a_k} = \infty}$.  We write this as $\boxed{a_k \ll b_k}$. 
  
  We also say that $a_k$ \uline{grows at the same rate as} $b_k$ if  $\boxed{\displaystyle \lim_{k \to \infty}
  \frac{a_k}{b_k} = 1}$. We write this as $\boxed{a_k \sim b_k}$, and often say that $a_k$ is \textit{asymptotic} to $b_k$.

  Here are some common functions of $k$, listed according to their growth
  rates:
  \begin{equation*} 
  r^k \ll C \ll \ln k \ll k^m \ll k^n \ll a^k \ll b^k \ll k! \ll k^k 
  \end{equation*}
  
 where $0 < m < n$, $0 < r < 1 < a < b$ and $C$ is a constant.
 
  Here are some more useful facts:
  \begin{itemize}[itemsep=0pt,topsep=0pt]
  \item A sum is asymptotic to its asymptotically largest term.
  \item If $b_k \gg a_k$, then $b_k \gg c a_k$ for any constant $c$.
  For instance, since $k^3 \gg k^2$, $k^3 \gg 50000 k^2$.
\end{itemize}
\end{framed}


\pspace{\newpage}




\begin{enumerate}[resume]
\item  Asymptotically simplify the following expressions as much as you
can.  (That is, find a simpler expression that is asymptotic to the
given one, and try to make your simpler expression as simple as
possible.)
\begin{enumerate}
  \item
  \begin{problem}[\vfill]
  $\displaystyle \frac{k^{1000} + k^{999} + 1.1^k}{\ln k + k^4}$.  
  \end{problem}

  \begin{solution}
  In the sum $k^{1000} + k^{999} + 1.1^k$ (in the numerator), $1.1^k$
  is asymptotically the largest of the three terms; that is, $1.1^k
  \gg k^{1000}$ and $1.1^k \gg k^{999}$.  So, $k^{1000} + k^{999} +
  1.1^k \sim 1.1^k$.

  In the sum $\ln k + k^4$ (in the denominator), $k^4 \gg \ln k$, so
  $\ln k + k^4 \sim k^4$.  Therefore, $\displaystyle \frac{k^{1000} +
  k^{999} + 1.1^k}{\ln k + k^4} \sim \boxed{\frac{1.1^k}{k^4}}$.
  \end{solution}

  \item
  \begin{problem}[\vfill]
  $\displaystyle \frac{k! + \pi k^{17} + 10000^k}{10^{10} \ln k + k^k 
  + k!}$.
  \end{problem}

  \begin{solution}
  In the sum $k! + \pi k^{17} + 10000^k$ (in the numerator), $k!$ is
  asymptotically the largest of the three terms.  In the sum $\ln k +
  k^k + k!$ (in the denominator), $k^k$ is asymptotically the largest
  of the three terms.  So, $\displaystyle \frac{k! + \pi k^{17} +
  10000^k}{\ln k + k^k + k!} \sim \frac{k!}{k^k}$. 
  \end{solution}

  \item
  \begin{problem}[\vfill]
  $\displaystyle \frac{2^k (k^3 + k^2)}{3^k}$.
  \end{problem}

  \begin{solution}
  In the sum $k^3 + k^2$, $k^3 \gg k^2$, so $k^3 + k^2 \sim k^3$.
  Therefore, $\displaystyle \frac{2^k (k^3 + k^2)}{3^k} \sim
  \boxed{\frac{2^k k^3}{3^k}}$.  (Since there are no longer any sums,
  we can't asymptotically simplify any further.)
  \end{solution}

  \item
  \begin{problem}[\vfill]
  $\displaystyle \frac{1}{k! (k^9 - 7 k^2)^{5/3}}$.
  \end{problem}

  \begin{solution}
  In the sum $k^9 - 7 k^2$, $k^9 \gg k^2$, so $k^9 - 7 k^2 \sim k^9$.
  Therefore, $\displaystyle \frac{1}{k! (k^9 - 7 k^2)^{5/3}} \sim
  \frac{1}{k! (k^9)^{5/3}} = \boxed{\frac{1}{k! \cdot k^{15}}}$.
  \end{solution}
  
  
  \item
  \begin{problem}[\vfill]
  $\displaystyle \frac{5 + 3 \sin^2 k}{k^3 + k!}$.
  \end{problem}

  \begin{solution}
  Since $k! \gg k^3$, $k^3 + k! \sim k!$, so $\displaystyle \frac{5 +
  3 \sin^2 k}{k^3 + k!} \sim \boxed{\frac{5 + 3 \sin^2 k}{k!}}$.  We
  can't asymptotically simplify the numerator, as neither of the
  summands grows faster than the other.
  \end{solution}

  \item
  \begin{problem}[\vfill]
  $\displaystyle \frac{(k + \sin k) \ln k}{5^k \cdot k^3}$.
  \end{problem}

  \begin{solution}
  $k + \sin k \sim k$, which we can check using the definition of
  $\sim$. Therefore, $\displaystyle \frac{(k + \sin k) \ln k}{5^k
  \cdot k^3} \sim \frac{k \ln k}{5^k \cdot k^3} = \boxed{\frac{\ln
  k}{5^k \cdot k^2}}$.
  \end{solution}
\end{enumerate}

\pspace{\newpage}

 \begin{framed}
\centerline{\textbf{Asymptotic Comparison Test}}
Suppose $\displaystyle \sum a_k$ and $\displaystyle \sum b_k$ are two series with \textbf{positive} terms (eventually) such that $a_k \sim b_k$. Then either both series converge or both series diverge.
\end{framed}





\item 
\note{Students often try to asymptotically simplify products in strange
ways.  Emphasize that we really only asymptotically simplify sums.}

Use asymptotics to help you decide whether the following series converge
or diverge.

\begin{enumerate}
  \item % from Spring 2007 midterm 2
  \begin{problem}[\vfill]
  $\displaystyle \sum_{k=1}^\infty \frac{5 \sqrt{k}}{k^2 + 2k}$.
  \end{problem}

  \begin{solution}
  $\displaystyle \frac{5 \sqrt{k}}{k^2 + 2k} \sim \frac{5
  \sqrt{k}}{k^2} = \frac{5}{k^{3/2}}$.  Since $\displaystyle
  \frac{5}{k^{3/2}} > 0$, we can apply the Asymptotic Comparison Test
  to conclude that $\displaystyle \sum \frac{5 \sqrt{k}}{k^2 + 2k}$
  behaves the same way (in terms of convergence) as $\displaystyle
  \sum \frac{5}{k^{3/2}}$.  The series $\displaystyle \sum
  \frac{5}{k^{3/2}}$ is a constant multiple of the $p$-series with $p
  = \frac{3}{2} > 1$, so it converges.  Therefore, $\displaystyle
  \sum_{k=1}^\infty \frac{5 \sqrt{k}}{k^2 + 2k}$ converges as well.
  \end{solution}
  
  
  \item
  \begin{problem}[\vfill]
  $\displaystyle \sum_{k=1}^\infty \frac{k^2}{(100 + 5k^2 +
  k^{11})^{1/3}}$
  \end{problem}

  \begin{solution}
  We know that $k^{11} \gg 100$ and $k^{11} \gg 5k^2$, so $100 + 5k^2
  + k^{11} \sim k^{11}$.  Thus, $\displaystyle \frac{k^2}{(100 + 5k^2
  + k^{11})^{1/3}} \sim \frac{k^2}{\left( k^{11} \right)^{1/3}} =
  \frac{k^2}{k^{11/3}} = \frac{1}{k^{5/3}}$.  Since $\displaystyle
  \frac{1}{k^{5/3}} > 0$, we can apply the Asymptotic Comparison Test
  to conclude that $\displaystyle \sum_{k=1}^\infty \frac{k^2}{(100 +
  5k^2 + k^{11})^{1/3}}$ behaves the same way (in terms of
  convergence) as $\displaystyle \sum \frac{1}{k^{5/3}}$.  This latter
  series is a $p$-series with $\displaystyle p = \frac{5}{3} > 1$, so
  it converges.  Therefore, $\displaystyle \sum_{k=1}^\infty
  \frac{k^2}{(100 + 5k^2 + k^{11})^{1/3}}$ converges as well.
  \end{solution}

  \item
  \begin{problem}[\vfill]
  $\displaystyle \sum_{k=1}^\infty \frac{k^4 + 2^k - 7 \ln k}{(3^k -
  10 k^2) k!}$.
  \end{problem}

  \begin{solution}
  $\displaystyle \frac{k^4 + 2^k - 7 \ln k}{(3^k - 10 k^2) k!} \sim
  \frac{2^k}{3^k \cdot k!} = \frac{(2 / 3)^k}{k!}$.  Since
  $\displaystyle \frac{(2 / 3)^k}{k!} > 0$, we can apply the
  Asymptotic Comparison Test to conclude that $\displaystyle \sum
  \frac{k^4 + 2^k}{(3^k - 10 k^2) k!}$ behaves the same way (in terms
  of convergence) as $\displaystyle \sum \frac{(2 / 3)^k}{k!}$.

  Now, we need to figure out whether $\displaystyle \sum \frac{(2 /
  3)^k}{k!}$ converges.  We can use the Direct Comparison Test:
  $\displaystyle 0 < \frac{(2 / 3)^k}{k!} \leq \left( \frac{2}{3}
  \right)^k$, and we know that $\displaystyle \sum \left( \frac{2}{3}
  \right)^k$ converges because it's a geometric series with common
  ratio $\displaystyle \frac{2}{3}$, and $\displaystyle \left|
  \frac{2}{3} \right| < 1$.  Therefore, the Comparison Test says
  that $\displaystyle \sum \frac{(2 / 3)^k}{k!}$ converges, so
  $\displaystyle \sum \frac{k^4 + 2^k - 7 \ln k}{(3^k - 10 k^2) k!}
  \sim \frac{2^k}{3^k \cdot k!}$ converges, too.
  \end{solution}
  
  
  \item
  \begin{problem}[1in]
  $\displaystyle \sum_{k=1}^{\infty} \frac{k^3 (e - 2 \cos
  k)}{\sqrt{k^7- 5 k}}$
  \end{problem}

  \begin{solution}
  First, here (in \ideacolor) is some intuitive reasoning that we
  might use to decide whether the series converges: \ideas{$e - 2 \cos
  k$ stays between $e - 2$ and $e + 2$, so it doesn't seem to be
  doing a lot in terms of how the terms go to 0.  The rest of the
  term, $\dfrac{k^3}{\sqrt{k^7 - 5k}}$, is asymptotic to
  $\displaystyle \frac{k^3}{\sqrt{k^7}} = \frac{1}{k^{1/2}}$, so I
  think the series should diverge because it behaves roughly like a
  $p$-series with $p = 1/2$.}

  To turn this into a mathematical argument, we need to use a
  convergence test.  We cannot use asymptotics to simplify $e - 2 \cos
  k$, so we won't be able to use the Asymptotic Comparison Test
  directly.  Instead, let's try using the Direct Comparison Test.
  Since we are trying to show that our series diverges, we need to
  show that its terms are \textit{bigger} than the terms of another
  series that we know to be divergent.  

  Since $-1 \leq \cos k \leq 1$, we have the inequality $e - 2 \leq e
  - 2 \cos k \leq e + 2$.  Therefore,
  \begin{align*}
  \frac{k^3 (e - 2 \cos k)}{\sqrt{k^7 - 5k}} & \geq \frac{k^3 (e -
  2)}{\sqrt{k^7 - 5k}} \textrm{ (right side has smaller numerator,
  same denominator)} \\
  & \geq \frac{k^3 (e - 2)}{\sqrt{k^7}} \textrm{ (same numerator,
  bigger denominator)} \\
  &= (e - 2) \frac{1}{k^{1/2}} \\
  & > 0 
  \end{align*}
  Since $\displaystyle (e - 2) \sum \frac{1}{k^{1/2}}$ diverges (it's
  a constant times a $p$-series with $p = 1/2$), we can conclude that
  the series in question also \fbox{diverges}, by the \textbf{Direct
  Comparison Test}.
  \end{solution}

\end{enumerate}



\end{enumerate}



\end{document}